% BHU PYQ -- Manually transcribed
% Paper: MTB-102 : Geometry
% Exam: B.Sc. (Hons.) Semester I Examination, 2023-24

\documentclass[11pt,a4paper]{article}
\usepackage[a4paper,top=2.5cm,bottom=2.5cm,left=2.8cm,right=2.8cm,headheight=14pt]{geometry}
\usepackage{amsmath,amssymb}
\usepackage[T1]{fontenc}
\usepackage[utf8]{inputenc}
\usepackage{lmodern}
\usepackage{microtype}
\usepackage{fancyhdr}
\usepackage{enumitem}
\usepackage{hyperref}
\usepackage{lastpage}
\usepackage{parskip}

\hypersetup{colorlinks=true,urlcolor=black,linkcolor=black,
  pdftitle={MTB-102 Geometry -- BHU B.Sc. Sem I 2023-24}}

\pagestyle{fancy}\fancyhf{}
\renewcommand{\headrulewidth}{0.4pt}\renewcommand{\footrulewidth}{0.4pt}
\fancyhead[L]{\small   \textit{Banaras Hindu University}}
\fancyhead[R]{\small   \textit{MTB-102: Geometry}}
\fancyfoot[C]{\small Page  \thepage\ of \pageref{LastPage}}


\newlist{parts}{enumerate}{2}
\setlist[parts,1]{label=(\alph*),leftmargin=*,itemsep=5pt,topsep=3pt}
\setlist[parts,2]{label=(\roman*),leftmargin=*,itemsep=3pt,topsep=2pt}

\newcommand{\pts}[1]{\hfill   \textit{[#1]}}


\begin{document}

\begin{center}
  {\large\bfseries BANARAS HINDU UNIVERSITY}\\[2pt]
  {\normalsize B.Sc.\ (Hons.) Science --- Semester I Examination, 2023-24}\\[6pt]
  \rule{0.8\linewidth}{0.5pt}\\[6pt]
  {\large\bfseries Paper No.\ MTB-102: Geometry}\\[4pt]
  {\normalsize Time: 3 Hours \hfill Maximum Marks: 70}\\[2pt]
  {\small Ref.\ No.:\ 067936}
\end{center}
\rule{\linewidth}{0.4pt}
\smallskip



\noindent   \textit{Instructions: Answer   \textbf{five} questions including Question~1 which is compulsory. Symbols have their usual meanings.}
\bigskip




\noindent   \textbf{Question 1.}   \textit{(Compulsory --- 2 marks each.)}

\begin{parts}
  \item Find the centre and radius of the circle $r = 3(\sqrt{3}\cos \theta + \sin \theta)$.
  \item Obtain the points on the conic $\dfrac{15}{r} = 3 + 5\cos \theta$ whose radius vector is $3$.
  \item Find the distance between two points $(2,\,\pi/3)$ and $(3,\,\pi/6)$ in polar coordinates.
  \item Find the angle between the planes $\mathbf{r}\cdot(3\hat{i}-4\hat{j}+5\hat{k})=3$ and $\mathbf{r}\cdot(2\hat{i}-\hat{j}-2\hat{k})=2$.
  \item Find the intercepts made on the axes by the plane $\mathbf{r}\cdot(2\hat{i}+3\hat{j}-2\hat{k})=6$.
  \item Write the line $\dfrac{x-1}{2}=y+1=\dfrac{z-3}{4}$ in vector form.
  \item Find the centre of the circle $x^2+y^2+z^2=6$, $z=2$.
  \item Define \emph{reciprocal cones}.
  \item Define \emph{right circular cylinder}.
  \item Find the equation of the normal to the conicoid $3x^2-y^2+2z^2=2$ at the point $(1,-3,2)$.
  \item The plane $x+2y-2z=4$ touches the paraboloid $3x^2+4y^2=24z$. Find the point of contact.
\end{parts}
\medskip\rule{\linewidth}{0.2pt}

\medskip


\noindent   \textbf{Question 2.}

\begin{parts}
  \item Define tangent at a point of a conic. Find the equation of the chord joining $(r_1,\alpha)$ and $(r_1,\beta)$ on the conic $\ell/r = 1+e\cos \theta$.\pts{6}
  \item Show that the director circle of the conic $\ell/r = 1+e\cos \theta$ is $\ell^2(1-e^2)+2e\ell r\cos \theta - 2r^2=0$.\pts{6}
\end{parts}
\medskip\rule{\linewidth}{0.2pt}

\medskip


\noindent   \textbf{Question 3.}

\begin{parts}
  \item A variable plane at constant distance $p$ from the origin meets the axes in $A,B,C$. Show the locus of the centroid of tetrahedron $OABC$ is $\dfrac{1}{x^2}+\dfrac{1}{y^2}+\dfrac{1}{z^2}=\dfrac{16}{p^2}$.\pts{6}
  \item Find the equation of the line through position vector $3\hat{i}+2\hat{j}-4\hat{k}$, parallel to both $\mathbf{r}\cdot(\hat{i}-2\hat{j}+3\hat{k})=5$ and $\mathbf{r}\cdot(3\hat{i}+\hat{j}-2\hat{k})=6$.\pts{6}
\end{parts}
\medskip\rule{\linewidth}{0.2pt}

\medskip


\noindent   \textbf{Question 4.}

\begin{parts}
  \item Find the length and equation of the shortest distance between $\mathbf{r}=\mathbf{a}+t\mathbf{b}$ and $\mathbf{r}=\mathbf{c}+s\mathbf{d}$.\pts{6}
  \item Find the equation of the sphere through $(1,1,0)$, $(1,0,1)$, $(0,1,1)$ with smallest radius.\pts{6}
\end{parts}
\medskip\rule{\linewidth}{0.2pt}

\medskip


\noindent   \textbf{Question 5.}

\begin{parts}
  \item Prove that $x+2y-z=4$ cuts $x^2+y^2+z^2-x+z-2=0$ in a circle of radius unity. Find the sphere with this circle as a great circle.\pts{6}
  \item Find the cone with vertex $(-2,4,3)$ and guiding curve $3x+2y=0$, $y^2+z^2=0$.\pts{6}
\end{parts}
\medskip\rule{\linewidth}{0.2pt}

\medskip


\noindent   \textbf{Question 6.}

\begin{parts}
  \item Find the enveloping cone of $x^2+y^2+z^2-3x+2y-4=0$ with vertex $(2,1,3)$.\pts{6}
  \item Find the right circular cylinder of radius $3$ with axis $\dfrac{x-2}{1}=\dfrac{y-1}{1}=\dfrac{z+3}{1}$.\pts{6}
\end{parts}
\medskip\rule{\linewidth}{0.2pt}

\medskip


\noindent   \textbf{Question 7.}

\begin{parts}
  \item Find the tangent planes of $2x^2-6y^2+3z^2=5$ passing through the line $x+9y-3z=0$, $3x-3y+6z=5$.\pts{6}
  \item Prove that enveloping cylinders of $\dfrac{x^2}{a^2}+\dfrac{y^2}{b^2}+\dfrac{z^2}{c^2}=1$ whose generators are parallel to $x=0$, $\dfrac{y}{b}=\dfrac{z}{c}$, meet $z=0$ in circles.\pts{6}
\end{parts}

\vfill\rule{\linewidth}{0.4pt}

\begin{center}\small
\href{https://bhu-exam.vercel.app/}{Click here to visit our website}\\[4pt]
\href{https://me-aryan.vercel.app/}{Click here to visit our contributor}\\[4pt]
\href{https://quashan-me.vercel.app/}{Click here to visit our contributor}
\end{center}
\end{document}
